\documentclass[11pt]{article}
\usepackage[utf8]{inputenc}\usepackage[T1]{fontenc}\usepackage[french]{babel}
\usepackage{amsmath,amssymb}\usepackage[margin=2.2cm]{geometry}
\begin{document}
\section*{Pourquoi $A$ et $\sigma$ sont anticorrélés dans un fit gaussien}

\paragraph{1. Covariance d'un fit par moindres carrés.}
La matrice de covariance des paramètres est l'inverse de la matrice de Fisher :
\begin{equation}
C = F^{-1}, \qquad
F_{ij} = \frac{1}{s^2}\sum_{\text{pixels}} \frac{\partial f}{\partial p_i}\,\frac{\partial f}{\partial p_j},
\end{equation}
où $s$ est le bruit sur un pixel et $p_i$ les paramètres du fit.

\paragraph{2. Dérivées d'une gaussienne.}
Pour $f(x) = A\,e^{-x^2/2\sigma^2}$ :
\begin{equation}
\frac{\partial f}{\partial A} = e^{-x^2/2\sigma^2} > 0,
\qquad
\frac{\partial f}{\partial \sigma} = A\,\frac{x^2}{\sigma^3}\,e^{-x^2/2\sigma^2} \geq 0 .
\end{equation}
Les deux dérivées ont le même signe sur tous les pixels, donc $F_{A\sigma} > 0$.

\paragraph{3. Inversion.}
Pour une matrice $2\times2$, l'élément hors diagonale de l'inverse change de signe :
\begin{equation}
C_{A\sigma} = -\frac{F_{A\sigma}}{\det F} < 0 .
\end{equation}

\paragraph{4. Valeur de la corrélation} (1D, sans offset, sommes remplacées par des intégrales) :
\begin{equation}
F_{AA} = \sqrt{\pi}\,\sigma,\qquad
F_{A\sigma} = \frac{\sqrt{\pi}}{2}\,A,\qquad
F_{\sigma\sigma} = \frac{3\sqrt{\pi}}{4}\,\frac{A^2}{\sigma},
\end{equation}
\begin{equation}
\rho(A,\sigma) = -\frac{F_{A\sigma}}{\sqrt{F_{AA}F_{\sigma\sigma}}} = -\frac{1}{\sqrt{3}} \simeq -0.58 .
\end{equation}
En 2D avec offset, la simulation donne $\rho(A,\sigma_x)\simeq\rho(A,\sigma_y)\simeq-0.49$ et $\rho(\sigma_x,\sigma_y)\simeq0.05$.

\paragraph{5. Conséquence sur le nombre d'atomes.}
Avec $N \propto A\,\sigma_x\,\sigma_y$, la propagation complète s'écrit
\begin{equation}
\left(\frac{\delta N}{N}\right)^2 =
\left(\frac{\delta A}{A}\right)^2 + \left(\frac{\delta \sigma_x}{\sigma_x}\right)^2 + \left(\frac{\delta \sigma_y}{\sigma_y}\right)^2
+ \frac{2\,\mathrm{Cov}(A,\sigma_x)}{A\,\sigma_x} + \frac{2\,\mathrm{Cov}(A,\sigma_y)}{A\,\sigma_y} + \frac{2\,\mathrm{Cov}(\sigma_x,\sigma_y)}{\sigma_x\,\sigma_y}.
\end{equation}
Les deux covariances avec $A$ sont négatives et $\mathrm{Cov}(\sigma_x,\sigma_y)\simeq0$ : ne garder que les trois premiers termes \emph{surestime} $\delta N/N$.
Simulation : $\delta N/N = 0.53\,\%$ (diagonale seule) contre $0.33\,\%$ (covariance complète).

\paragraph{Intuition.} Si le fit élargit un peu la gaussienne, il baisse l'amplitude pour garder la même intégrale : le nombre d'atomes est mieux déterminé que $A$ et $\sigma$ séparément.
\end{document}
